\documentclass[12pt]{article}
\PassOptionsToPackage{table}{xcolor}
\usepackage[utf8]{inputenc}
\usepackage[T1]{fontenc}
\usepackage[italian]{babel}
\usepackage{multicol}
\usepackage{enumitem}
\usepackage{array}
\usepackage{xcolor}
\usepackage{graphicx}
\usepackage{tikz}
\usetikzlibrary{
    shapes.geometric,
    arrows.meta,
    positioning,
    calc,
    patterns,
    decorations.pathmorphing
}
\usepackage[margin=2cm]{geometry}
\usepackage{amsmath,amssymb}
\setlength{\parindent}{0pt}
\definecolor{headercolor}{RGB}{220,220,220}
\newcounter{quesito}
\setcounter{quesito}{0}
\newcommand{\nuovoquesito}{\stepcounter{quesito}\textbf{Domanda \arabic{quesito}.}}
\newcommand{\itemdomanda}[2]{%
    \par\noindent
    \begin{minipage}[t]{\linewidth}
        \nuovoquesito \\ #1
        \begin{enumerate}[label=(\Alph*), nosep, leftmargin=*]
            #2
        \end{enumerate}
    \end{minipage}
    \vspace{10pt}
}
\begin{document}
\section*{Simulazione}
\hfill Data: 16 apr 2026

\itemdomanda{Tutte le rose sono fiori. Alcuni fiori appassiscono presto. Quindi:}{
    \item Tutte le rose appassiscono presto
    \item Alcune rose appassiscono presto
    \item Nessuna rosa appassisce presto
    \item È possibile che alcune rose appassiscano presto
    \item I fiori che appassiscono sono solo rose
}

\itemdomanda{Date le seguenti premesse:
    \begin{enumerate}
        \item Tutti i canarini ben nutriti cantano a squarciagola.
        \item Nessun canarino che canti a squarciagola è malinconico.
    \end{enumerate}
    Quale conclusione è logicamente corretta?}{
    \item Nessun canarino ben nutrito è malinconico
    \item Qualche canarino ben nutrito è malinconico
    \item Tutti i canarini malinconici sono ben nutriti
    \item I canarini che non cantano sono sicuramente ben nutriti
    \item Chi canta a squarciagola è malinconico
}

\itemdomanda{La condizione di esistenza di $x^{-1} + y^{-1}$ è
}{
\item $ x \neq 1 \lor y \neq 1 $
\item $ x \neq -1 \lor y \neq -1 $
\item tutti i reali
\item $ x \neq y $
\item $ x \neq 0 \land y \neq 0 $
}

\itemdomanda{Qual è il resto della divisione del polinomio $P(x) = x^{3} - 4x + 1$ per $(x - 2)$?}{
\item $1$
\item $0$
\item $-1$
\item $17$
\item $3$
}

\itemdomanda{Se $ 4^{x-1} = 1 $, allora $ x $ vale:
}{
\item $ 4 $
\item $ 0 $
\item $ 1 $
\item $ 2 $
\item $ -1 $
}

\itemdomanda{$ \log_{10}(0.01) $ vale:
}{
\item $ 2 $
\item $ 1 $
\item $ -2 $
\item $ -1 $
\item $ 0.1 $
}

\end{document}
